Retain AR as the reference generator and treat range relabeling as a
conditioning/reweighting policy. Its gain does not establish that another
acquired batch is beneficial, or that LOW requests improve as much as HIGH
requests. Future comparisons should preserve a common assessment history
and separate fixed-bin fidelity from extension of the observed range.

The next latent test should target information that survives posterior
sampling and produces a directed change in generated structures. Recovering
ordering from tiny encoder-mean variations is useful diagnosis, but does not
establish that the learned latent improves generation. The entropy study
likewise shows why the sampling objective matters: its 1,000\,K proposal
helps that target but loses to uniform sampling at 5,000\,K.

\begin{itemize}[leftmargin=*,itemsep=3pt]
\item \textbf{Ricardo} reviews endpoint structures, symmetry and units,
interprets teacher disagreements, and defines later DFT/relaxation checks.
\item \textbf{Beihao} reproduces AR/CVAE training and develops a controlled
sampled-latent test with the existing data, matched budgets, and genuine,
condition-only, shuffled and random-encoder comparisons.
\item \textbf{Varad} reproduces the paired evaluations, audits thresholds
and denominators, reports per-cell variability and negative results, and
maintains the analysis and cost receipts.
\end{itemize}

Independent DFT accuracy, relaxed structures, non-O3 phases, size transfer,
and stable dynamics remain physical validation tasks. The current
fixed-lattice MLIP results do not establish those outcomes.
